%
% system.tex
%
% Lux Technical Architecture Specification
% Written for principal engineers evaluating the system.
%
% Compile: pdflatex system.tex
%

\documentclass[a4paper,10pt]{article}
\usepackage[margin=1in]{geometry}
\usepackage{booktabs}
\usepackage{listings}
\usepackage{xcolor}
\usepackage{enumitem}
\usepackage{longtable}
\usepackage[colorlinks=true,linkcolor=blue,citecolor=blue,urlcolor=blue]{hyperref}

\lstset{
  basicstyle=\ttfamily\small,
  keywordstyle=\color{blue!80!black},
  commentstyle=\color{green!50!black},
  stringstyle=\color{red!60!black},
  breaklines=true,
  frame=single,
  framerule=0.3pt,
  rulecolor=\color{gray!50},
  backgroundcolor=\color{gray!5},
  columns=fullflexible,
}

\newcommand{\modpath}[1]{\path{src/punt_lux/#1}}
\newcommand{\tool}[1]{\texttt{#1}}
\newcommand{\msg}[1]{\texttt{#1}}
\newcommand{\elem}[1]{\texttt{#1}}
\newcommand{\field}[1]{\texttt{#1}}

\begin{document}

\title{Lux: Current Architecture Specification}
\author{Visual Output Surface for AI Agents --- a System in Transition}
\date{July 2026 --- v0.21.0}
\hypersetup{
  pdftitle={Lux: Current Architecture Specification},
  pdfauthor={Punt Labs},
  pdfsubject={Architecture},
  pdfcreator={pdflatex}
}
\maketitle

\tableofcontents
\newpage

% ===================================================================
\section{Overview}
% ===================================================================

Lux is a visual output surface for Claude Code agents and local
applications.  An agent or an app describes a user interface as a tree of
typed JSON elements and submits it to Lux; Lux installs that interface in
one authoritative store and renders it in a native ImGui window.

The system is built from one engine fronted by thin client surfaces.  The
engine is split into two cooperating processes:

\begin{itemize}
  \item The \textbf{Hub} (\tool{luxd}) owns the authoritative interface
    state, runs the real interaction handlers, and coordinates the display.
    It is the single front door: every client talks to the Hub.
  \item The \textbf{Display} (\tool{lux-display}) holds a full replica of
    the interface and renders it every frame.  It captures user input and
    forwards each interaction back to the Hub.  It never talks to a client
    directly.
\end{itemize}

Four client surfaces reach the Hub, and each is thin:

\begin{itemize}
  \item The \textbf{library} surface is the operations layer itself ---
    a direct in-process call into the engine.
  \item The \textbf{MCP} surface is a set of tools an agent calls; it
    speaks streamable HTTP to the Hub.
  \item The \textbf{REST} surface is a typed FastAPI application on the
    Hub; it is the front door for the command-line tool and any non-MCP
    caller.
  \item The \textbf{CLI} (\tool{lux}) is an HTTP client of the REST
    surface.
\end{itemize}

\subsection{Current and Target}

This document describes the architecture \emph{as the code stands today}.
The canonical design intent lives in
\texttt{docs/architecture/target/target.md}; where this document and that
target disagree, the target defines the \emph{direction} and this document
reports the \emph{present}.  Each place the present diverges from the
target is called out explicitly.

The code has converged substantially on the target.  The Hub is
authoritative, the operations layer is the single home of every
capability, the MCP tools and REST routes are adapters over it, and the
command-line tool reaches the Hub over REST.  Two areas remain in
transition:

\begin{itemize}
  \item \textbf{Element migration.}  Fifteen of the twenty-five element
    kinds are on the Element-ABC path (data and behavior on the class);
    ten kinds are still frozen dataclasses decoded by per-kind codecs.
    See Section \S\ref{sec:hubdisplay}.
  \item \textbf{The display's own state.}  The Display still owns its
    render loop, its themes and window settings, and its diagnostic ring
    buffers.  The Hub reaches those facts by proxying a read or write over
    its one connection to the Display (Section \S\ref{sec:tools}).
\end{itemize}

\subsection{Companion Documents}

The sibling files provide narrower supporting views:

\begin{itemize}
  \item \texttt{target/target.md} --- the canonical rewrite target.  Start
    there for design intent; return here for current behavior.
  \item \texttt{one-code-path.md} --- the design of the operations layer,
    the REST front door, the thin adapters, and the streamable-HTTP
    transport.  It is the design record for most of what this document
    now describes as shipped (ADR DES-055).
  \item \texttt{target/introspection-api.md} --- the verification and
    control surface an agent uses to inspect live Hub and Display state.
\end{itemize}

Four formal models hold parts of the system to a checked specification;
they are cross-referenced where each applies and listed together in
Section \S\ref{sec:formal}.

\subsection{Design Principles}

\begin{enumerate}[label=\textbf{P\arabic*}.]
  \item \textbf{One engine, thin clients.}  Every capability is one typed
    operation on the Hub.  The library call, the MCP tool, the REST route,
    and the CLI command all run that same operation; none reimplements it.
  \item \textbf{The Hub is the authority.}  Interface state that a client
    submits is owned by the Hub.  The Display is a replica, and it loses
    every disagreement.
  \item \textbf{Protocol is the API surface.}  An agent describes an
    interface as JSON; Lux renders it.  There is no shared memory and no
    required Python import.
  \item \textbf{UI state crosses the wire; render calls do not.}
    Serialized elements, scene updates, and interaction messages travel
    between Hub and Display.  ImGui calls stay inside the Display.
  \item \textbf{Whole-UI resend on change.}  When Hub-side state changes,
    the Hub re-sends the whole affected interface and the Display replaces
    its copy.  There is no incremental diff protocol.
\end{enumerate}

% ===================================================================
\section{System Architecture}
% ===================================================================

\subsection{Component Map}

The package divides into the engine and its client surfaces.  The engine
is the operations layer, the Hub domain, the scene state, and the
protocol; the surfaces are the MCP tools, the REST routers, the CLI, and
the transport plumbing that carries each.  The Display is the engine's
second process.

\begin{longtable}{p{2.6cm}p{3.3cm}p{8.2cm}}
\toprule
\textbf{Concern} & \textbf{File or package} & \textbf{Responsibility} \\
\midrule
\endfirsthead
\toprule
\textbf{Concern} & \textbf{File or package} & \textbf{Responsibility} \\
\midrule
\endhead
Operations    & \modpath{operations/}            & The engine's front of
                house.  One \texttt{Operations} facade composes concern
                classes --- \texttt{SceneOperations},
                \texttt{QueryOperations}, \texttt{MenuOperations},
                \texttt{DisplayControlOperations}, \texttt{PubSubOperations},
                \texttt{ConvenienceOperations}, \texttt{DisplayModeOperations}.
                Each method takes a typed request and returns a typed
                result or an \texttt{OpError}.  Models live in
                \modpath{operations/models/}. \\
Hub domain    & \modpath{domain/hub/}            & The authoritative
                engine core.  \texttt{HubDisplay} is the store;
                \texttt{HubReplicator} is the sole writer to the Display;
                \texttt{FrameLifecycle}, \texttt{FrameExpiry}, and
                \texttt{ExpirySweep} own the frame time-to-live;
                \texttt{HubReads} answers introspection;
                \texttt{HubMenuRegistry} owns the menu bar;
                \texttt{HubSceneWriter} applies updates and clears;
                \texttt{ClientRegistry} owns luxd's one display connection
                and the D21 dispatch. \\
Domain (pure) & \modpath{domain/}                & The renderer-free and
                socket-free element model.  \texttt{element\_abc.py} is
                the \texttt{Element} ABC (data and behavior);
                \texttt{interaction.py} and \texttt{container\_interaction.py}
                are the typed events; \texttt{handlers/remote\_dispatch.py}
                wraps handlers for the two-tier path. \\
Scene         & \modpath{scene/}                 & The Display's scene
                graph.  \texttt{SceneManager} owns the unframed scenes and
                per-scene widget state; \texttt{FrameBook} owns the frames
                and the scene-to-frame and scene-to-owner maps;
                \texttt{frame.py} is the \texttt{Frame} window. \\
Protocol      & \modpath{protocol/}              & The wire contract:
                length-prefixed JSON framing, the \texttt{FrameReader}
                streaming buffer, the twenty-five element kinds in
                \modpath{protocol/elements/}, and the message families in
                \modpath{protocol/messages/}. \\
Display       & \modpath{display/}               & The ImGui render loop
                and coordinator.  \texttt{server.py},
                \texttt{element\_renderer.py}, \texttt{table\_renderer.py},
                \texttt{menu\_manager.py}, \texttt{idle\_screen.py}, and
                \texttt{texture\_cache.py}. \\
MCP tools     & \modpath{tools/}                 & Thin adapters over the
                operations facade.  \texttt{tools.py} and
                \texttt{composite\_tools.py} are display-facing tools;
                \texttt{subscribe\_tools.py} are the session pub-sub tools;
                \texttt{server.py} owns the \texttt{FastMCP} instance. \\
REST          & \modpath{rest/}                  & The typed FastAPI
                surface.  \texttt{RestSurface} composes concern routers
                (\texttt{scenes}, \texttt{menus}, \texttt{display},
                \texttt{config}); \texttt{status.py} is the one
                \texttt{OpError}-to-HTTP table. \\
MCP transport & \modpath{mcp\_transport.py}      & luxd's MCP leg:
                streamable HTTP mounted at \texttt{/mcp} beside the REST
                routes.  \texttt{mcp\_endpoint.py} resolves session
                identity; \texttt{mcp\_session.py} runs the per-session
                cleanup cascade; \texttt{session\_cleanup.py} tears down a
                departed session's Hub state. \\
CLI transport & \modpath{rest\_client.py}        & \texttt{LuxRestClient}
                --- the CLI's typed HTTP client of luxd's REST routes.
                \texttt{rest\_transport.py} is the HTTP transport;
                \texttt{rest\_loopback.py} is the in-process transport
                tests use. \\
Gateway       & \modpath{luxd.py}                & Builds the one FastAPI
                app on uvicorn, mounts the MCP leg and the REST surface,
                and starts the frame-expiry sweep on the event loop. \\
Display client & \modpath{display\_client.py}    & \texttt{DisplayClient}
                --- the Unix-socket client of the Display.  Constructed in
                exactly one place, \modpath{domain/hub/clients.py}, and an
                AST guard enforces that (Section \S\ref{sec:ipc}). \\
Paths         & \modpath{paths.py}               & \texttt{DisplayPaths}
                --- display socket discovery, PID lifecycle, and
                auto-spawn.  \modpath{hub\_paths.py} is the symmetric
                \texttt{HubPaths} for luxd's port and PID files. \\
Apps/Beads    & \modpath{apps/beads.py}          & \texttt{BeadsBrowser}
                --- builds the beads issue table as an element tree.  Pure:
                it holds no socket. \\
Config        & \modpath{config.py}              & \texttt{ConfigManager}
                --- reads and writes the per-repo display-mode toggle in
                \texttt{<repo>/.punt-labs/lux.md}. \\
CLI           & \modpath{\_\_main\_\_.py}        & The Typer app: show,
                service, status, and installation commands. \\
\bottomrule
\end{longtable}

\subsection{The Operations Facade}\label{sec:operations}

Every capability the system offers is one method on the
\texttt{Operations} facade (\modpath{operations/facade.py}).  The facade
composes seven concern classes, each owning one group of capabilities, and
delegates each method to the right one.  A caller --- an MCP adapter, a
REST route, or a test --- holds one object and calls one method.

The facade binds no process singletons at import time.  Its
\texttt{for\_store} classmethod receives every collaborator (the
\texttt{HubDisplay} store, the replicator, the Hub, the client registry,
the menu registry, the display connection) from the composition root, so
the operations are testable without the running process.

\subsubsection{Typed Requests and the Never-Raising Parse}

Each operation takes a typed request model and returns either the
operation's own success model or a shared error.  Validation is the
operation's job, not the adapter's.  A request model carries a
never-raising \texttt{parse} classmethod: it validates raw arguments and
returns either the request or an \texttt{OpError}, so an adapter that
builds a plain mapping and calls \texttt{parse} can never raise past its
string contract.  The operation accepts the parse result and passes an
\texttt{OpError} straight through.

The consequence is one schema per capability, applied at one point.  A bad
\field{layout} becomes an \texttt{invalid\_request} error inside
\texttt{parse}; the MCP adapter renders it to the legacy status string, and
a REST route renders it to an HTTP 422.  The rules live in one place; only
the rendering differs by surface.

\subsubsection{The Error Contract}

Every operation may return \texttt{OpError}
(\modpath{operations/models/common.py}), a frozen model tagged
\texttt{kind="error"} so a result cannot be both a success and a failure.
The \field{code} is a closed set the caller branches on; the \field{reason}
is the human sentence for logs and messages.

\begin{center}
\begin{tabular}{lp{9cm}}
\toprule
\textbf{Code} & \textbf{Meaning} \\
\midrule
\texttt{display\_unavailable} & The display process is not running. \\
\texttt{timeout}             & A proxied round-trip to the display
  exceeded its bound. \\
\texttt{rejected}            & The engine refused the caller's request
  (the caller's fault). \\
\texttt{invalid\_request}    & The request itself did not type-check. \\
\texttt{not\_found}          & The named scene or resource does not exist. \\
\texttt{fault}               & An engine-side failure: a malformed display
  reply, an unreadable config file. \\
\bottomrule
\end{tabular}
\end{center}

The REST surface maps each code to one HTTP status through a single shared
table (\modpath{rest/status.py}): \texttt{invalid\_request} to 422,
\texttt{not\_found} to 404, \texttt{rejected} to 409, \texttt{fault} to
502, \texttt{display\_unavailable} to 503, and \texttt{timeout} to 504.
The mapping is a table, not per-route logic, so a new operation inherits it
for free.  \texttt{LuxRestClient} reads the same table in reverse, so the
CLI observes the same contract from the other end.

\subsection{Data Flow}

Three data paths connect a client to the Display.

\subsubsection{Write Path (Client $\to$ Hub $\to$ Display)}

\begin{enumerate}
  \item A client calls a capability --- the \tool{show} MCP tool, a
    \texttt{PUT /scenes/\{id\}} REST call, or \tool{lux show beads}.
  \item The surface builds a plain mapping and calls the request builder,
    then calls one operation.  For \tool{show} that is
    \texttt{Operations.render}.
  \item The operation validates the request, then installs the scene into
    \texttt{HubDisplay} --- the authoritative store --- under the store
    lock, indexing every element and its owner.
  \item The operation marks the scene dirty on the \texttt{HubReplicator}
    and returns its typed result at once.  The client's call does not wait
    on the Display.
  \item The \texttt{HubReplicator}, the one background worker that writes
    to the Display, drains the dirty set, serializes the scene, and sends
    it over luxd's single socket connection.
  \item The Display receives the scene, replaces its replica, and renders
    it every frame.
\end{enumerate}

\subsubsection{Event Path (User $\to$ Hub, two-tier D21)}

The event path is where the Hub/Display split is clearest.  The Display
does not run interaction handlers; it forwards each interaction to the Hub,
where the real handler runs.

\begin{enumerate}
  \item When the Display installs a scene, it wraps every handler on each
    element, replacing each event bucket with a
    \texttt{RemoteDispatchGroup} that captures the socket-send closure.
  \item The user interacts with a widget.  The renderer fires the typed
    event on the Display's element copy, and the wrapped handler runs
    instead of the real handler.
  \item The wrapper builds one \msg{RemoteEventHandlerInvocation}
    (\field{element\_id}, \field{action}, \field{event\_kind},
    \field{value}, \field{ts}, \field{scene\_id}) and sends it to the Hub.
    One interaction produces one message per element/event bucket, however
    many handlers that bucket holds.
  \item The Hub receives the message in \texttt{ClientRegistry.\_hub\_interaction\_dispatch},
    resolves the authoritative element by \field{(scene\_id, element\_id)},
    checks its owner, builds the typed event, and fires the real handler
    chain once on its own copy.
  \item If the handler mutated Hub state (for example a dialog dismissed
    itself), the Hub marks the scene dirty; the replicator re-sends the
    whole scene and the Display renders the current tree.
\end{enumerate}

A \field{frame\_close} action is not an element interaction: the Display
sends it when the user closes a frame, and the Hub answers by removing that
frame's scenes so both tiers agree the frame is gone (Section
\S\ref{sec:frames}).

\subsubsection{Update Path (Incremental Patching)}

\begin{enumerate}
  \item A client calls \tool{update(scene\_id, patches)}, reaching
    \texttt{Operations.update}.
  \item The operation applies the patch batch to \texttt{HubDisplay}
    through \texttt{HubSceneWriter}: a \field{set} patch changes fields on
    the target element, a \field{remove} patch drops it.
  \item The operation marks the scene dirty; the replicator re-sends the
    whole scene, and the Display renders the result.
\end{enumerate}

% ===================================================================
\section{The Hub/Display Slice}\label{sec:hubdisplay}
% ===================================================================

This section describes the Hub-authoritative core and the Element-ABC path
the migration is converging on.  It is where new element work lands and the
reference for how the whole system is meant to behave once migration
completes.

The domain code in \modpath{domain/} imports no ImGui, no socket, and no
JSON.  Adapters wrap it: the Display renders it, and the replicator moves
its serialized elements over the wire.

\subsection{HubDisplay: the Authoritative Store}

\texttt{HubDisplay} (\modpath{domain/hub/hub\_display.py}) is the Hub's
single authoritative surface for scene state.  It is a facade over typed
collaborators, each with one responsibility:

\begin{center}
\begin{tabular}{lp{8.5cm}}
\toprule
\textbf{Collaborator} & \textbf{Responsibility} \\
\midrule
\texttt{ElementIndex}       & \field{(scene\_id, element\_id)} $\to$
  \texttt{Element} lookup. \\
\texttt{OwnerTracker}       & Every element's owning connection; the
  interaction path checks it, and the disconnect path walks it. \\
\texttt{RootRegistry}       & Scene-root elements that take part in the
  property-observer cascade. \\
\texttt{ChildIndex}         & Parent-to-children edges captured at install
  time so cascade removal drops every descendant in one walk. \\
\texttt{HubClientRegistry}  & The connections currently registered as Hub
  clients. \\
\texttt{FrameLifecycle}     & Each scene's frame presentation, its optional
  TTL deadline, and its teardown --- all under the store lock (Section
  \S\ref{sec:frames}). \\
\texttt{HubReads}           & The authoritative introspection reads that
  answer \tool{inspect\_scene} and \tool{list\_scenes}. \\
\bottomrule
\end{tabular}
\end{center}

A write dispatches on the typed \texttt{Update} sum ---
\texttt{AddElement}, \texttt{SetProperty}, \texttt{RemoveElement} --- and
enforces ownership: a connection cannot mutate or evict another
connection's elements.  \texttt{AddElement} recurses into composite
children so a button buried inside a dialog is indexed and later
click-resolvable.  Re-showing a scene removes the connection's previous
roots and installs the new ones through the same write path, so ownership,
root observers, and child indexes rebuild in one place.

\subsection{The Element ABC}

The migrated kinds subclass the \texttt{Element} ABC in
\modpath{domain/element\_abc.py}.  The ABC carries data \emph{and}
behavior:

\begin{itemize}
  \item \textbf{Composite render template.}  \texttt{render()} is a
    template method and is never overridden; a composite overrides
    \texttt{\_children()} to return its children, and a leaf inherits the
    empty default.
  \item \textbf{Handler registry.}  \texttt{add\_handler},
    \texttt{remove\_handler}, and \texttt{fire} key handlers by
    \texttt{Event} subclass.  \texttt{fire} snapshots the bucket so a
    handler that mutates the registry mid-dispatch cannot affect the
    in-flight call.
  \item \textbf{Property-observer surface.}  \texttt{add\_observer},
    \texttt{mark\_removed}, and \texttt{removed} let a parent composite
    react to a child being removed.  All three removal paths --- an agent
    \texttt{RemoveElement}, a component self-dismiss, and a connection
    disconnect --- route through the one \texttt{mark\_removed} method.
  \item \textbf{Native serialization.}  \texttt{\_\_reduce\_\_} and
    \texttt{\_\_setstate\_\_} let an element cross the Hub-to-Display
    boundary.  Hub-only bookkeeping (the observer closures that reference
    \texttt{HubDisplay}) is dropped on serialization; the Display is a
    replica and does not need Hub observers.  Handlers survive so the
    Display can wrap them for remote dispatch.
\end{itemize}

\subsection{Two-Tier Handler Dispatch (D21)}

The Hub says: ``I am sending you a copy of an interface.  If anyone
interacts with it, I will do the work.''  The Display says: ``I will render
that copy and forward interactions back to the Hub.''

Both tiers hold the same element state and both call
\texttt{elem.fire(event)}.  The difference is what the handler does:

\begin{itemize}
  \item \textbf{Hub:} the real handler runs --- \texttt{call\_model},
    \texttt{publish}, \texttt{mark\_removed}.
  \item \textbf{Display:} the wrapped handler sends one
    \msg{RemoteEventHandlerInvocation} to the Hub, which resolves the
    authoritative element, checks its owner, and fires the real event once.
\end{itemize}

Four event kinds cross this boundary today:

\begin{center}
\begin{tabular}{llp{6.5cm}}
\toprule
\textbf{Event} & \textbf{Wire kind} & \textbf{Fired by} \\
\midrule
\texttt{ButtonClicked}  & \texttt{button\_clicked}  & a button press. \\
\texttt{ValueChanged}   & \texttt{value\_changed}   & a checkbox, slider,
  combo, radio, selectable, or text-input change. \\
\texttt{HeaderToggled}  & \texttt{header\_toggled}  & a
  \texttt{collapsing\_header} opening or closing. \\
\texttt{TabChanged}     & \texttt{tab\_changed}     & a \texttt{tab\_bar}
  switching tab. \\
\bottomrule
\end{tabular}
\end{center}

\texttt{ButtonClicked} and \texttt{ValueChanged} are in
\modpath{domain/interaction.py}; \texttt{HeaderToggled} and
\texttt{TabChanged} are in \modpath{domain/container\_interaction.py}.  The
Hub builds the typed event from the wire kind and value, denying by default
when the value has the wrong shape.  Grouping is per element/event bucket:
a button with three handlers still produces one remote message, and the Hub
replays the full chain once.

\subsection{Agent Subscribe / Publish}

The Hub also owns an application-level publish/subscribe channel
(\modpath{domain/hub/hub.py}), distinct from the element-observer mechanism
above; the two share no machinery.

Every operation is scoped to a connection.  A connection cannot see, touch,
or publish into another connection's topics.  \texttt{subscribe} registers
the caller's outbound writer against a topic; \texttt{publish} fans an
\msg{ObserverMessage} payload out to that connection's subscribers using
snapshot-then-iterate, so a slow handler cannot stall concurrent publishes.
Topics --- \texttt{openTicket}, \texttt{closeTicket},
\texttt{markTicketInProgress} --- are defined by the app or agent, not by
Lux.  This is the business-event channel, not the interface-replication
path.

\subsection{Migrated vs.\ Still-Legacy}

\begin{center}
\begin{tabular}{lp{9.5cm}}
\toprule
\textbf{State} & \textbf{What it covers} \\
\midrule
\textbf{Migrated}     & The fifteen Element-ABC kinds: \texttt{text},
  \texttt{button}, \texttt{checkbox}, \texttt{dialog}, \texttt{progress},
  \texttt{input\_text}, \texttt{input\_number}, \texttt{slider},
  \texttt{color\_picker}, \texttt{combo}, \texttt{radio},
  \texttt{selectable}, \texttt{group}, \texttt{collapsing\_header}, and
  \texttt{tab\_bar}.  Three of these --- \texttt{group},
  \texttt{collapsing\_header}, \texttt{tab\_bar} --- are ABC containers,
  migrated only when their whole subtree is ABC.  \texttt{HubDisplay}
  authority with ownership and cascade removal, the D21 four-event path,
  the frame lifecycle, and Agent Subscribe/Publish are all in this
  column. \\
\textbf{Transitional} & The Display process, which runs the D21 event path
  yet still owns rendering, scene tabs, frames, and its own diagnostic
  buffers as in the legacy design. \\
\textbf{Legacy}       & The ten not-yet-migrated element kinds as frozen
  dataclasses: \texttt{markdown}, \texttt{image}, \texttt{separator},
  \texttt{spinner}, \texttt{draw}, \texttt{plot}, \texttt{table},
  \texttt{tree}, \texttt{window}, and \texttt{modal}. \\
\bottomrule
\end{tabular}
\end{center}

The direction is fixed by \texttt{target/target.md}; the pace is set by
migrating one element family at a time, container plus primitives first
(the migration plan, DES-041).

% ===================================================================
\section{Wire Protocol}
% ===================================================================

\subsection{Framing}

Every message is a length-prefixed JSON frame:

\begin{lstlisting}
+------------------+----------------------------+
| 4 bytes (uint32) |  N bytes (UTF-8 JSON)      |
| big-endian       |  payload                   |
+------------------+----------------------------+
\end{lstlisting}

Maximum payload size is 16\,MiB ($2^{24}$ bytes).  The
\texttt{FrameReader} (in \modpath{protocol/\_\_init\_\_.py}) accumulates
partial reads into a byte buffer and yields complete frames via
\texttt{drain()} (raw dicts) or \texttt{drain\_typed()} (Message
dataclasses).  An oversize frame raises \texttt{ValueError}, which the
socket layer converts to a client disconnect.

This framing is the Hub-to-Display wire.  A client reaches the Hub over
HTTP (MCP streamable HTTP, or REST), not over this socket; the socket is
Hub-internal (Section \S\ref{sec:ipc}).

\subsection{Message Types}

The message layer is split across six families in
\modpath{protocol/messages/}: \texttt{scene}, \texttt{lifecycle},
\texttt{menu}, \texttt{introspect}, \texttt{observer}, and
\texttt{remote\_invocation}.  There is no standalone \texttt{interaction}
family --- a user interaction travels as \msg{RemoteEventHandlerInvocation}
from the \texttt{remote\_invocation} module.  A \texttt{MessageRegistry}
(\modpath{protocol/messages/registry.py}) is the wire dispatch table,
populated at import time by each family module; a test can build an isolated
registry without touching the module default.

\subsubsection{Client $\to$ Display}

\begin{center}
\begin{tabular}{lp{9cm}}
\toprule
\textbf{Type} & \textbf{Purpose} \\
\midrule
\msg{SceneMessage}  & Replace or add a scene.  Fields: \field{id},
  \field{elements[]}, \field{title}, \field{layout}, \field{frame\_id},
  \field{frame\_title}, \field{frame\_size}, \field{frame\_flags},
  \field{frame\_layout}. \\
\msg{UpdateMessage} & Incremental patches.  Fields: \field{scene\_id},
  \field{patches[]} (each: \field{id}, \field{set}, \field{remove}). \\
\msg{ClearMessage}  & Remove scenes and reset state. \\
\msg{PingMessage}   & Heartbeat.  Field: \field{ts}. \\
\msg{ConnectMessage} & Client identity handshake.  Field: \field{name}. \\
\msg{MenuMessage}   & Set the menu bar. \\
\msg{RegisterMenuMessage} & Register per-connection menu items. \\
\msg{ThemeMessage}  & Apply a theme by name. \\
\msg{QueryRequest}  & A control or introspection request the display
  answers.  Fields: \field{method}, \field{params}. \\
\bottomrule
\end{tabular}
\end{center}

\subsubsection{Display $\to$ Client}

\begin{center}
\begin{tabular}{lp{9cm}}
\toprule
\textbf{Type} & \textbf{Purpose} \\
\midrule
\msg{ReadyMessage}       & Handshake on connect: protocol version and
  capabilities. \\
\msg{AckMessage}         & Acknowledges a scene or update: \field{scene\_id},
  \field{ts}, optional \field{error}. \\
\msg{RemoteEventHandlerInvocation} & A user interaction forwarded to the
  owning Hub: \field{element\_id}, \field{action}, \field{event\_kind}
  (\texttt{button\_clicked}, \texttt{value\_changed},
  \texttt{header\_toggled}, \texttt{tab\_changed}), \field{value},
  \field{ts}, optional \field{scene\_id}. \\
\msg{ObserverMessage}    & A connection-scoped business event for the Agent
  Subscribe/Publish path: \field{topic} + \field{payload}. \\
\msg{PongMessage}        & Ping response. \\
\msg{QueryResponse}      & A control or introspection response:
  \field{method}, \field{result}, optional \field{error}. \\
\bottomrule
\end{tabular}
\end{center}

An unknown message type deserializes as \msg{UnknownMessage} for forward
compatibility rather than disconnecting the peer.

\subsection{Element Kinds}

The wire layer exposes twenty-five element kinds plus the non-element
\texttt{Patch} payload used by \msg{UpdateMessage}.  Fifteen kinds are on
the Element-ABC path (Section \S\ref{sec:hubdisplay}); the remaining ten
are frozen, slotted dataclasses registered into \texttt{ElementCodec}.
The \texttt{Element} union is assembled in
\modpath{protocol/elements/\_\_init\_\_.py}.

An ABC kind serializes two ways.  Its Level-1 form is a plain JSON dict
through its per-kind codec; its wire form crossing the Hub-to-Display
boundary is a base64-encoded pickled entry inside the \msg{SceneMessage}
codec.  A legacy kind crosses as a plain dict either way.

\begin{longtable}{llp{6cm}}
\toprule
\textbf{Category} & \textbf{Element} & \textbf{Notes} \\
\midrule
\endfirsthead
\toprule
\textbf{Category} & \textbf{Element} & \textbf{Notes} \\
\midrule
\endhead
\texttt{basics}    & \elem{text}              & Content plus style
              (heading, caption, success, error, code); optional
              \field{color} hex string.  ABC. \\
            & \elem{markdown}          & CommonMark via imgui\_md.  Legacy. \\
            & \elem{image}             & File path or base64 data.  Legacy. \\
            & \elem{separator}         & Visual divider.  Legacy. \\
            & \elem{progress}          & Fraction bar with label.  ABC. \\
            & \elem{spinner}           & Loading indicator.  Legacy. \\
\midrule
\texttt{inputs}    & \elem{button}            & Label, action, disabled;
              \field{arrow} and \field{small} variants.  ABC. \\
            & \elem{slider}            & Min, max, value, integer mode.
              ABC. \\
            & \elem{checkbox}          & Boolean toggle; fires
              \texttt{ValueChanged}.  ABC. \\
            & \elem{combo}             & Dropdown selection.  ABC. \\
            & \elem{input\_text}       & Single-line text input.  ABC. \\
            & \elem{radio}             & Radio button group.  ABC. \\
            & \elem{input\_number}     & Numeric input with step buttons and
              clamping.  ABC. \\
            & \elem{color\_picker}     & RGB/RGBA hex picker.  ABC. \\
            & \elem{selectable}        & Click-to-select item; fires
              \texttt{ValueChanged}.  ABC. \\
\midrule
\texttt{dialog}    & \elem{dialog}            & Dialog root with title,
              children, handlers, and a model-backed lifecycle.  ABC. \\
\midrule
\texttt{layout}    & \elem{group}             & Row, column, or paged
              layout.  ABC container. \\
            & \elem{collapsing\_header} & Expandable section; fires
              \texttt{HeaderToggled}.  ABC container. \\
            & \elem{tab\_bar}          & Tabbed container; fires
              \texttt{TabChanged}.  ABC container. \\
            & \elem{window}            & Movable, resizable floating panel.
              Legacy. \\
            & \elem{tree}              & Recursive node hierarchy;
              \field{flat} flag.  Legacy. \\
            & \elem{modal}             & Modal popup; emits \texttt{"closed"}.
              Legacy. \\
\midrule
\texttt{graphics}  & \elem{draw}              & Canvas with typed draw
              commands (\S\ref{sec:draw-commands}).  Legacy. \\
            & \elem{plot}              & ImPlot: line, scatter, bar series.
              Legacy. \\
\midrule
\texttt{table}     & \elem{table}             & Columns, rows, built-in
              filters, detail panel, pagination.  Legacy. \\
\bottomrule
\end{longtable}

\subsection{Draw Commands}\label{sec:draw-commands}

The \elem{draw} element holds a tuple of typed \texttt{DrawCommand}
records.  Each command is a frozen, slotted dataclass that satisfies the
\texttt{DrawCommand} Protocol (\texttt{TYPE: ClassVar[str]},
\texttt{to\_wire(self) -> dict}, \texttt{from\_wire(cls, d, *, ctx) ->
Self}).

\begin{center}
\begin{tabular}{llp{6cm}}
\toprule
\textbf{Module} & \textbf{Record} & \textbf{Geometry} \\
\midrule
\texttt{draw\_commands\_line}   & \texttt{Line}        & Start, end, color, thickness. \\
                                & \texttt{Polyline}    & Point list, color, thickness, closed flag. \\
\midrule
\texttt{draw\_commands\_shape}  & \texttt{Rect}        & Top-left, bottom-right, color, thickness, rounding, filled. \\
                                & \texttt{Circle}      & Center, radius, color, thickness, filled. \\
                                & \texttt{Triangle}    & Three points, color, thickness, filled. \\
\midrule
\texttt{draw\_commands\_curve}  & \texttt{BezierCubic} & Four control points, color, thickness. \\
\midrule
\texttt{draw\_commands\_text}   & \texttt{TextGlyph}   & Position, content, color. \\
\bottomrule
\end{tabular}
\end{center}

Shared value classes live in \texttt{draw\_values.py} and
\texttt{draw\_bounds.py}: \texttt{Point2} (a validated $[x, y]$ pair),
\texttt{Color} (\texttt{\#RRGGBB} / \texttt{\#RRGGBBAA} hex),
\texttt{Thickness} (a strictly positive stroke width), \texttt{Radius}
(strictly positive), and \texttt{Rounding} (non-negative).  The
\texttt{WireContext} in \texttt{draw\_wire.py} carries the per-error prefix
and the require-or-raise helpers.  \texttt{DrawCommandDecoder} in
\texttt{draw\_decoder.py} dispatches on the wire \texttt{cmd} string to the
right \texttt{from\_wire} classmethod; \texttt{DrawElement.commands} is
typed \texttt{tuple[DrawCommand, ...]} and the renderer dispatches via a
typed \texttt{match} statement.

% ===================================================================
\section{Scene and Frame Lifecycle}\label{sec:frames}
% ===================================================================

A scene targets a named \emph{frame} --- a persistent ImGui window that
organizes the workspace.  The lifecycle answers three questions the same
way on both tiers: when a frame appears, when it disappears, and how long
it lives without a refresh.

\subsection{Frames Die With Their Scenes}

A frame exists to show scenes.  When a frame's last scene is removed, the
frame closes; a frame still holding other scenes stays open.

The Hub blanks a scene by pushing it with no elements.  The Display treats
an empty-element push as a removal, not as an empty window: an unframed
empty push dismisses the scene, and a framed one dismisses the scene from
its frame and closes the frame once it holds nothing.  This is the amended
contract the legacy display-server model now checks (Section
\S\ref{sec:formal}).

When the user closes a frame on the Display, the Display sends a
\field{frame\_close} action.  The Hub removes every scene that frame held
from \texttt{HubDisplay} and marks each dirty, so the authoritative store
agrees with the window the user already closed.  The blank the replicator
then sends is idempotent with that close.

\subsection{Scenes Survive Their Session}

A scene outlives the connection that created it.  When an MCP session ends
--- the client disconnects, or the SDK reaps it on idle --- the Hub forgets
that connection as a client and releases its subscriptions and inbox, but
it leaves the scenes installed.  The scenes stay owned by the departed
connection id, so a later explicit removal --- a frame close, a
\tool{clear}, or a TTL expiry --- can still reference the owner.  Nothing
is blanked, so no repaint is marked.  A session's death does not erase what
it drew.

\subsection{Frame Time-to-Live}

A frame may carry a time-to-live.  Shown with a \field{ttl\_seconds}, it is
removed once that many seconds pass with no re-show refreshing it.  Three
small components implement this, all in \modpath{domain/hub/}:

\begin{itemize}
  \item \texttt{FrameExpiry} owns the bookkeeping: a \field{frame\_id} to a
    deadline on a monotonic clock (a duration, immune to wall-clock
    adjustment).  It answers which frames are due now and how long until
    the soonest one.  It holds no lock.
  \item \texttt{FrameLifecycle} owns each scene's frame presentation, its
    deadline, and its teardown, all under the one store lock.  Recording a
    presentation and arming its deadline is a single locked step, so a
    re-show that refreshes a frame cannot install it yet forget to re-arm
    it.
  \item \texttt{ExpirySweep} is the task on luxd's event loop that enforces
    the deadlines.  It sleeps until the soonest deadline (capped at a
    one-second idle poll), sweeps the frames now due, and marks their
    scenes dirty so the replicator blanks both tiers --- the same path a
    manual frame close takes.  It is one more store mutator on the loop,
    not a second timer thread.
\end{itemize}

Because arming a deadline and expiring a frame both run under the store
lock, a frame re-armed with a fresh TTL is never torn down by a stale
deadline: the re-show and the sweep serialize, and whichever runs second
sees the other's effect.  A re-show carrying no TTL makes the frame
permanent.  This atomicity is the safety property the frame-expiry formal
model checks (Section \S\ref{sec:formal}).

A due frame is peeked, not consumed: the sweep tears each frame down and
only then disarms its deadline, so a tear-down that raises leaves the frame
armed to retry on the next sweep rather than stranded --- torn down on the
Hub but never blanked on the Display.  Frames are isolated from one
another: a frame whose tear-down raises is logged and skipped, and every
frame that did tear down is still repainted.

\subsection{The Frame Window}

The \texttt{Frame} class (\modpath{scene/frame.py}) is the Display's window
for a set of scenes.  \texttt{FrameBook} (\modpath{scene/frame\_book.py})
owns the frame collection and the scene-to-frame and scene-to-owner maps,
split out of \texttt{SceneManager} so each has one job.

\begin{center}
\begin{tabular}{lp{8cm}}
\toprule
\textbf{Field} & \textbf{Type and Purpose} \\
\midrule
\field{frame\_id}       & \texttt{str} --- unique identifier. \\
\field{title}           & \texttt{str} --- window title bar. \\
\field{scenes}          & \texttt{dict[str, SceneMessage]} --- the frame's
  scenes, keyed by \field{scene\_id}. \\
\field{scene\_order}    & \texttt{list[str]} --- tab ordering. \\
\field{active\_tab}     & \texttt{str | None} --- the visible scene. \\
\field{minimized}       & \texttt{bool} --- docked to the bottom bar. \\
\field{cascade\_index}  & \texttt{int} --- stagger offset for initial
  placement. \\
\field{initial\_size}   & \texttt{tuple[int,int] | None} --- first-use size
  hint. \\
\field{flags}           & ImGui window flags (\texttt{no\_resize},
  \texttt{no\_collapse}, \texttt{no\_title\_bar}, \texttt{no\_background},
  \texttt{no\_scrollbar}, \texttt{auto\_resize}). \\
\field{layout}          & \texttt{Literal["tab", "stack"]} --- multi-scene
  composition mode. \\
\bottomrule
\end{tabular}
\end{center}

Left-clicking the docking background opens the World menu: list frames,
restore minimized frames, close all frames, and reach display settings.  A
minimized frame appears as a button in a bottom dock bar; clicking it
restores the frame.

% ===================================================================
\section{Transport and Process}\label{sec:ipc}
% ===================================================================

luxd runs one FastAPI application on one uvicorn port (default 8430).  Two
surfaces mount on that app: the MCP leg at \texttt{/mcp}, and the typed
REST routers beside it.  A \texttt{/health} route reports liveness and the
live session count.

\subsection{The MCP Leg: Streamable HTTP}

luxd's MCP surface is streamable HTTP (\modpath{mcp\_transport.py}), the
MCP SDK's current transport, mounted as an ASGI application at \texttt{/mcp}
on the same FastAPI app as the REST routes.  There is no separate WebSocket
endpoint, and mcp-proxy no longer sits in lux's path: Claude Code connects
to \texttt{http://127.0.0.1:8430/mcp} directly through its native HTTP MCP
configuration.

Each MCP session has its own identity.  \texttt{McpAsgiApp}
(\modpath{mcp\_endpoint.py}) resolves the session identity on each request
and refuses the reserved REST connection key.  \texttt{SessionScopedServer}
(\modpath{mcp\_session.py}) runs the Hub cleanup cascade when a session
ends, and \texttt{session\_cleanup.py} tears down that session's Hub-side
scenes stay-behind, menu items, subscriptions, and inbox.  A vanished
client --- a killed process, a dropped socket --- is reaped 30 minutes
after its last activity, so its Hub state cannot leak until restart.
Session teardown is isolated: one session ending does not disturb another.

\subsection{The Loopback Policy}

luxd binds loopback by default and refuses a non-loopback \texttt{--host}
at startup with a clear message, so the bind and the trust policy can never
disagree.  One \texttt{LoopbackTransportPolicy}
(\modpath{transport\_policy.py}) is the single source of the loopback
allowlist: it decides which hosts luxd starts on, and it projects the same
allowlist into the SDK's transport security settings so the streamable-HTTP
transport rejects a foreign \texttt{Host} header (DNS-rebinding) or
\texttt{Origin} header (cross-site).  Off-loopback access needs
authentication that is not built yet; the multi-machine future in
\texttt{target.md} arrives together with a bearer token and a
bind-derived origin policy.

\subsection{The Display Socket Is Hub-Internal}

The Display listens on a Unix domain socket, and luxd is its only client.
Nothing else in the package opens that socket.  \texttt{DisplayClient}
(\modpath{display\_client.py}) is the class that opens it, and it is
constructed in exactly one place: \texttt{ClientRegistry} in
\modpath{domain/hub/clients.py}, which holds luxd's one lazy connection and
its reconnect policy.

An AST guard enforces this
(\path{tests/domain/test_display_socket_internal.py}).  The guard parses
every module in the package and fails if any module outside the allowed set
--- \modpath{display\_client.py} and \modpath{domain/hub/clients.py} ---
imports or references \texttt{DisplayClient}.  Detection is by parse tree,
not by regex, so every spelling (an aliased import, a multi-line
parenthesized import, a whole-module import, a bare or attribute reference)
is caught at once, while a mention in a docstring or comment is never
mistaken for use.  A new reach-around cannot slip back in through a module
an enumerated list would have forgotten.

\subsection{Display Socket Discovery and Auto-Spawn}

The Display's socket path resolves in priority order (\modpath{paths.py}):

\begin{enumerate}
  \item \texttt{\$LUX\_SOCKET}.
  \item \texttt{\$XDG\_RUNTIME\_DIR/lux/display.sock}.
  \item \texttt{/tmp/lux-\$USER/display.sock} (fallback, chosen to stay
    within macOS's $\sim$104-character \texttt{AF\_UNIX} path limit).
\end{enumerate}

A PID file is written at \texttt{\{socket\_path\}.pid} on startup, and luxd
verifies liveness via \texttt{os.kill(pid, 0)} before connecting.
\texttt{DisplayPaths.ensure()} spawns the Display if it is not running,
detached from the caller's process group, and polls socket existence and
PID liveness every 100\,ms until ready or a 5\,s timeout.  When a send to
the Display fails with an \texttt{OSError}, the replicator drops the stale
client so the next send reconnects.

\subsection{The Display Render Loop}

The Display runs an ImGui render loop via \texttt{hello\_imgui.run()}.  Each
frame, the \texttt{DisplayServer.\_on\_frame()} callback runs four
non-blocking phases: accept new socket connections, poll the connection for
messages and dispatch each by type, render the active scenes and frames,
and flush any queued interaction messages back to the Hub.  A frame
completes in microseconds when idle; \texttt{hello\_imgui.fps\_idling.fps\_idle
= 30.0} caps the idle rate, and GLFW vsync paces active frames.

Table filtering runs in the render loop with zero Hub round-trips.
\texttt{TableRenderer} draws the search and combo filter controls, applies
AND logic across active filters (case-insensitive substring for search,
exact match for combos), resets pagination to page zero on a filter change,
and renders a two-column detail grid for the selected row.

% ===================================================================
\section{Client Surfaces}\label{sec:tools}
% ===================================================================

Every surface is a thin adapter over the operations facade
(\S\ref{sec:operations}).  An MCP tool parses its arguments, calls one
operation, and formats the result; it holds no validation, no scene
construction, and no display round-trips.

\subsection{MCP Tools}

The display-facing tools live in \modpath{tools/tools.py} and
\modpath{tools/composite\_tools.py}; the session pub-sub tools live in
\modpath{tools/subscribe\_tools.py}.

\begin{longtable}{lp{8.2cm}}
\toprule
\textbf{Tool} & \textbf{Operation and purpose} \\
\midrule
\endfirsthead
\toprule
\textbf{Tool} & \textbf{Operation and purpose} \\
\midrule
\endhead
\multicolumn{2}{l}{\textbf{Scenes}} \\
\midrule
\tool{show}           & \texttt{render} --- install a scene from element
  dicts. \\
\tool{show\_table}    & \texttt{render\_table} --- compose a filterable
  table tree and delegate to \texttt{render}. \\
\tool{show\_dashboard} & \texttt{render\_dashboard} --- compose a metrics,
  charts, and table tree and delegate to \texttt{render}. \\
\tool{update}         & \texttt{update} --- apply patches by element ID. \\
\tool{clear}          & \texttt{clear} --- remove the caller's scenes. \\
\midrule
\multicolumn{2}{l}{\textbf{Menus}} \\
\midrule
\tool{set\_menu}      & \texttt{set\_menu} --- replace the Hub-owned menu
  bar. \\
\tool{register\_tool} & \texttt{register\_menu\_item} --- add a tool item
  for the caller's session. \\
\tool{list\_menus}    & \texttt{list\_menus} --- read the menu bar. \\
\midrule
\multicolumn{2}{l}{\textbf{Introspection (Hub-authoritative)}} \\
\midrule
\tool{inspect\_scene} & \texttt{inspect\_scene} --- a scene's element tree
  from the store, with each element's \field{render\_path}
  (\texttt{"abc"} or \texttt{"legacy"}) and \field{resolved\_props}. \\
\tool{list\_scenes}   & \texttt{list\_scenes} --- every live scene and
  frame from the store. \\
\tool{list\_clients}  & \texttt{list\_clients} --- the Hub's sessions and
  their scopes. \\
\midrule
\multicolumn{2}{l}{\textbf{Display facts (proxied)}} \\
\midrule
\tool{get\_display\_info} & \texttt{get\_display\_info} --- backend,
  geometry, FPS, PID, uptime. \\
\tool{get\_window\_settings} & \texttt{get\_window\_settings} --- opacity,
  font scale, decoration, idle FPS. \\
\tool{set\_window\_settings} & \texttt{set\_window\_settings}. \\
\tool{get\_theme}     & \texttt{get\_theme} --- the active theme and the
  available themes. \\
\tool{set\_theme}     & \texttt{set\_theme}. \\
\tool{set\_frame\_state} & \texttt{set\_frame\_state} --- minimize or
  restore a frame. \\
\tool{list\_recent\_events} & \texttt{list\_recent\_events} --- the
  display's recent interactions. \\
\tool{list\_errors}   & \texttt{list\_errors} --- the display's recent
  errors. \\
\tool{screenshot}     & \texttt{screenshot} --- refuses with
  \texttt{rejected}: framebuffer capture is unsolved below the message
  layer (DES-028). \\
\tool{ping}           & \texttt{ping} --- display liveness probe. \\
\midrule
\multicolumn{2}{l}{\textbf{Per-repo config}} \\
\midrule
\tool{display\_mode}  & \texttt{read\_display\_mode}. \\
\tool{set\_display\_mode} & \texttt{write\_display\_mode}. \\
\midrule
\multicolumn{2}{l}{\textbf{Session pub-sub}} \\
\midrule
\tool{subscribe}      & \texttt{subscribe} --- subscribe the session to a
  topic in its own scope. \\
\tool{unsubscribe}    & \texttt{unsubscribe}. \\
\tool{publish}        & \texttt{publish} --- fan a business event to the
  session's subscribers. \\
\tool{recv}           & \texttt{receive} --- take the next session-scoped
  business event.  This is not the display interaction stream. \\
\bottomrule
\end{longtable}

An introspection tool reads Hub-authoritative state: \tool{inspect\_scene},
\tool{list\_scenes}, and \tool{list\_clients} ask the Hub, not the Display.
A display-fact tool --- a theme, a window setting, a framebuffer, the
display's own ring buffers --- proxies a read or write over luxd's one
connection.  The Hub cannot own an ImGui theme or a GPU backend string, so
those stay display facts; the caller still reaches them the same way,
through a Hub operation.  Menu state is the deliberate exception: a menu is
interface the agent submitted, so the Hub owns it and the replicator pushes
it like any scene.

\subsection{REST Routes}

The REST surface (\modpath{rest/}) binds each route body to a request
model, calls one operation, and returns the result model; FastAPI derives
the OpenAPI schema from the model, so there is no second schema to maintain.

\begin{center}
\begin{tabular}{llp{4.5cm}}
\toprule
\textbf{Method and path} & \textbf{Operation} & \textbf{Result} \\
\midrule
\texttt{PUT /scenes/\{id\}}          & \texttt{render}       & \texttt{SceneShown | OpError} \\
\texttt{PATCH /scenes/\{id\}}        & \texttt{update}       & \texttt{SceneShown | OpError} \\
\texttt{DELETE /scenes}              & \texttt{clear}        & \texttt{Cleared} \\
\texttt{GET /scenes}                 & \texttt{list\_scenes} & \texttt{SceneList} \\
\texttt{GET /scenes/\{id\}}          & \texttt{inspect\_scene} & \texttt{SceneInspection | OpError} \\
\texttt{GET /clients}                & \texttt{list\_clients} & \texttt{ClientList} \\
\texttt{GET /menus}                  & \texttt{list\_menus}  & \texttt{MenuList} \\
\texttt{PUT /menus}                  & \texttt{set\_menu}    & \texttt{Ok | OpError} \\
\texttt{POST /menus/items}           & \texttt{register\_menu\_item} & \texttt{Ok | OpError} \\
\texttt{GET /events}                 & \texttt{list\_recent\_events} & \texttt{RecentEvents} \\
\texttt{GET /errors}                 & \texttt{list\_errors} & \texttt{RecentErrors} \\
\texttt{GET /display}                & \texttt{get\_display\_info} & \texttt{DisplayInfo | OpError} \\
\texttt{GET,PUT /display/theme}      & \texttt{get/set\_theme} & \texttt{ThemeState | OpError} \\
\texttt{GET,PATCH /display/window}   & \texttt{get/set\_window\_settings} & \texttt{WindowSettings | OpError} \\
\texttt{PATCH /display/frames/\{id\}} & \texttt{set\_frame\_state} & \texttt{Ok | OpError} \\
\texttt{GET /display/ping}           & \texttt{ping}         & \texttt{Pong | OpError} \\
\texttt{GET,PUT /display-mode}       & \texttt{read/write\_display\_mode} & \texttt{DisplayModeState | OpError} \\
\bottomrule
\end{tabular}
\end{center}

Pub-sub is not on REST.  Publishing is scoped to the publisher's own
session, and a REST caller shares one anonymous default scope with no
subscribe route, so a REST publish could return success while being
structurally unable to reach any subscriber.  The pub-sub operations stay
MCP-only until a REST caller has a real session identity to subscribe under.
For the same reason, every REST call today lands in one default Hub scope;
per-caller REST scoping arrives with the multi-user future.

\subsection{The Command-Line Tool}

The CLI is the third thin client of the one engine.  \tool{lux show beads}
builds the beads element tree with the pure \texttt{BeadsBrowser} and sends
it to luxd's REST API through \texttt{LuxRestClient}
(\modpath{rest\_client.py}); it never touches the display socket.  The
client locates luxd's port from \texttt{HubPaths}, posts the request model,
and reads the typed result.  Two outcomes are distinct: luxd being
unreachable (no port file, a refused connection, a stalled response) raises
\texttt{HubUnavailableError} with an actionable message, while the Hub's
refusal of a reachable request comes back as a typed \texttt{OpError} mapped
from the HTTP status.

% ===================================================================
\section{Performance}
% ===================================================================

\subsection{Frame Budget}

\begin{center}
\begin{tabular}{lrl}
\toprule
\textbf{Parameter} & \textbf{Value} & \textbf{Notes} \\
\midrule
Idle FPS            & 30             & \texttt{fps\_idling.fps\_idle = 30.0} \\
Active FPS          & 60+            & GLFW vsync, GPU-bound \\
Socket poll timeout & 0              & Non-blocking \texttt{select()} \\
Max message size    & 16\,MiB        & Hard limit in \texttt{FrameReader} \\
Connection timeout  & 5.0\,s         & \texttt{ensure()} and handshake \\
Session idle reap   & 1800\,s        & MCP session idle timeout \\
Frame-expiry poll   & 1.0\,s         & \texttt{ExpirySweep} idle cap \\
Texture loading     & Lazy           & First render triggers PIL + OpenGL upload \\
Event ring buffer   & 200 entries    & Display recent interactions \\
Error ring buffer   & 100 entries    & Display recent errors \\
\bottomrule
\end{tabular}
\end{center}

\subsection{Critical Path}

A write does not wait on the Display.  An operation returns as soon as it
has updated \texttt{HubDisplay} and marked the scene dirty; the replicator
does the socket send on its own thread, so a stuck or dead Display never
blocks an agent-facing call.  This liveness property --- no mutator is ever
disabled by display I/O --- is one of the invariants the replicator formal
model checks (Section \S\ref{sec:formal}).

The Display's per-frame cost is a \texttt{select()} on the one connection,
a buffer drain, the ImGui draw calls (linear in visible element count), and
an event flush.  Table filtering is $O(R \cdot C)$ per frame in rows and
columns, fast for tables under $\sim$10{,}000 rows, with pagination
bounding the rendered rows.

\subsection{Memory}

\begin{itemize}
  \item \textbf{Scene storage.}  Scenes are Python dataclass trees in
    memory, bounded per message by the 16\,MiB cap; dismissed scenes are
    garbage-collected.
  \item \textbf{Texture cache.}  OpenGL texture IDs are held until exit
    with no eviction policy, so image-heavy sessions accumulate GPU memory.
  \item \textbf{Ring buffers.}  The event and error buffers are bounded
    deques; constant memory.
\end{itemize}

\subsection{Known Limitations}

\begin{enumerate}
  \item \textbf{No ImGui ID scoping.}  Widget IDs derive from
    \field{element\_id} alone, so ImGui internal state can bleed between
    tabs with colliding IDs.
  \item \textbf{No texture eviction.}  Long sessions with many images
    accumulate GPU memory.
  \item \textbf{Per-frame filtering.}  Table filters re-run every frame
    even when nothing changed; harmless at current scale.
  \item \textbf{No screenshot capture.}  Framebuffer capture is unsolved
    below the message layer, so \tool{screenshot} refuses (DES-028).
\end{enumerate}

% ===================================================================
\section{Security}
% ===================================================================

\subsection{Trust Boundaries}

There are two.  The first is luxd's HTTP surface: it binds loopback, refuses
a non-loopback bind at startup, and rejects a foreign \texttt{Host} or
\texttt{Origin} header (Section \S\ref{sec:ipc}).  The second is the
Display's Unix socket: luxd is its only client, filesystem permissions
protect the socket directory (user-only, mode \texttt{0700}), and an
oversize or malformed frame disconnects the peer rather than crashing the
Display.

\subsection{Render Function Consent}

A render function is the one element kind that executes arbitrary code, and
its security model is consent-based, not sandbox-based.  A best-effort AST
scan flags suspicious imports and dangerous builtins as a UX signal, not a
boundary; a consent modal shows the full source with any warnings, and the
user must click Allow to run it; execution is not sandboxed --- the code
runs in the Display's process.  A render function should be treated as
user-approved code, equivalent to a script the user chose to run.

\subsection{Protocol Robustness}

\begin{itemize}
  \item Malformed JSON disconnects the peer; it does not crash the Display.
  \item An unknown message type deserializes as \msg{UnknownMessage} for
    forward compatibility.
  \item An oversize frame is rejected by \texttt{FrameReader} and converted
    to a disconnect.
  \item A patch that tries to change an element's \field{id} or
    \field{kind} is dropped.
  \item Double-remove of a client is idempotent.
\end{itemize}

% ===================================================================
\section{Platform Integration}
% ===================================================================

\subsection{macOS}

\begin{itemize}
  \item The Display runs as a normal Dock app; the Dock icon and Cmd-Tab
    entry aid operational debugging.
  \item \texttt{setproctitle("Lux")} names the process in \texttt{ps} and
    \texttt{top}.
  \item Font discovery prefers Arial Unicode or Helvetica and merges Apple
    Symbols and STIX Two Math for mathematical glyphs.
  \item Socket paths stay within the $\sim$104-character \texttt{AF\_UNIX}
    limit; tests use \texttt{tempfile.mkdtemp(prefix="lux-")}.
\end{itemize}

\subsection{Linux}

\begin{itemize}
  \item Font discovery uses DejaVu Sans or Noto Sans and merges Noto Sans
    Symbols2 and Noto Sans Math.
  \item Socket placement prefers \texttt{\$XDG\_RUNTIME\_DIR/lux/}.
  \item Window managers use the X11 window title (``Lux'').
\end{itemize}

% ===================================================================
\section{Formal Models}\label{sec:formal}
% ===================================================================

Four Z specifications hold parts of the system to a checked model.  Each is
type-checked with Spivey's \texttt{fuzz} and model-checked with ProB; the
first two model concurrency the current code runs, and the third models the
Display singleton lifecycle.  The fourth is the legacy display-server model,
which the current display code is still refined against.

\begin{center}
\begin{tabular}{p{4.3cm}p{9cm}}
\toprule
\textbf{Model} & \textbf{What it proves} \\
\midrule
\texttt{docs/hub\_replicator.tex} & The single-writer, two-lock replication
  discipline.  Six invariants I1--I6: only the replicator writes to the
  Display (single writer); no mutator is ever blocked by display I/O
  (agent-path liveness); no thread holds the store lock and the send lock
  at once (no lock cycle); the display shows the store's latest state at
  rest, including after a reap and respawn (no lost update); a respawn
  follows a kill (one display); and the menu bar is repainted before rest
  (menu durability). \\
\texttt{docs/frame\_expiry.tex} & The re-show / expiry atomicity invariant:
  a frame re-armed with a fresh TTL is never torn down by a stale deadline,
  because both the re-show and the sweep run under the one store lock.  The
  fidelity variant drops the atomicity and ProB reaches the exact bad state
  where a freshly re-shown frame has been torn down (Section
  \S\ref{sec:frames}). \\
\texttt{docs/display\_lifecycle.tex} & The Display-singleton spawn
  lifecycle (DES-037/038): at most one Display exists across concurrent
  spawn attempts, and the lifecycle is deadlock-free. \\
\texttt{docs/architecture/ display-server.tex} & The legacy display-server
  state machine: client connect and disconnect, scene replacement,
  patching, and event flush.  Amended so that an empty-element push is a
  scene removal, matching the frames-die-with-scenes behavior (Section
  \S\ref{sec:frames}); the refinement harness now drives an empty push and
  holds the code to the amended model. \\
\bottomrule
\end{tabular}
\end{center}

Refinement tests in \texttt{tests/test\_display\_refinement.py} tie the
concrete Python back to the legacy model through an abstraction function,
and partition tests in \texttt{tests/test\_display\_partition.py} derive
cases from its operation schemas.  The legacy model predates multi-scene
tabs and the Hub/Display split; its abstraction maps multi-scene state to
the active tab.

% ===================================================================
\section{Test Architecture}
% ===================================================================

The suite is layered; the exact count changes, so this records the shape.

\begin{itemize}
  \item \textbf{Protocol and codec.}  Framing, serialization roundtrips,
    typed draw-command decode, value validation, and self-validation
    (DES-039).  Every element kind has a Level-1 roundtrip; a protocol
    change without one is unshippable.
  \item \textbf{Operations and surfaces.}  The operations facade over
    fakes, the REST routers via \texttt{TestClient}, the MCP adapters, and
    the CLI's \texttt{LuxRestClient}.
  \item \textbf{Hub domain.}  \texttt{HubDisplay} authority and ownership,
    the D21 dispatch, the frame lifecycle, and the replicator.
  \item \textbf{Guards.}  The display-socket AST guard
    (\S\ref{sec:ipc}) and the MCP-tool characterization corpus, which
    replays a captured snapshot per tool so the string contract cannot
    drift silently.
  \item \textbf{Formal-model alignment.}  The refinement and partition
    tests against the legacy display-server model.
  \item \textbf{Integration and e2e.}  Marker-gated tests exercise the real
    Hub/Display boundary; the business-event-loop harness
    (\texttt{tests/e2e/}) proves the full bidirectional D21 loop for a
    composed surface across the faithful in-memory connection.
\end{itemize}

Quality gates run through the Makefile: \texttt{ruff} (lint and format),
\texttt{mypy} and \texttt{pyright} (both strict), \texttt{pytest},
\texttt{check-oo} (the OO ratchet against \texttt{.oo-baseline.json}), and
the snapshot-parity gate.  \texttt{make check} is the composite gate.

\end{document}
